\label{exp:cea01}

\subsubsection*{Постановка}
\textbf{Физический вопрос.}
На~$M$ аннигиляция $p^+p^-$ --- локальное событие на решётке; наблюдатель на~$T$ видит осреднённое $|\Phi|$.
Появляется ли после pm на NN \emph{крупномасштабная} двухфронтовая структура: минимум и ось диполя в $\Delta|\Phi|$ на mid-plane readout?

\textbf{Что проверяем.}
\begin{enumerate}[label=\textbf{(\alph*)},leftmargin=*,itemsep=0.25em]
\item Сценарий \texttt{pm\_nn}: после аннигиляции coarse-поле $|\Phi|$ даёт $\ge 2$ выраженных пика и согласованную ось диполя.
\item Контроль \texttt{pp\_nn}: без аннигиляции --- иная картина (нет ложного диполя pm).
\item Связь микро-канала CE-F-01 с макро-прокси для последующего T-ансамбля CE-A-02.
\end{enumerate}

\textbf{Рабочая гипотеза.}
Осреднение на mid-plane после pm выделяет дипольную компоненту $\Delta|\Phi|$, отсутствующую в pp-контроле.

\textbf{Наблюдаемые.}
Число пиков $|\Phi|$, ориентация оси диполя, контраст pm vs pp.

\textbf{Критерий согласия с теорией.}
pm $\to$ дипольная прокси; pp $\to$ контроль без той же структуры; согласованность с линией CE-A-02.

\subsubsection*{Теоретическая часть}
Мост $M\to T$ здесь не спектроскопия: это проверка, что \emph{правильный} микро-канал оставляет след в readout, пригодный для статистики на~$T$.
Осреднение выполняется после эволюции pm, как в протоколе аннигиляционной серии.

\subsubsection*{Ход эксперимента}
Два NN-сценария: \texttt{pm\_nn}, \texttt{pp\_nn}; осреднение $|\Phi|$ на mid-plane.
Реестр: \texttt{Floor1\_pm\_coarse\_dipole}.

\begin{enumerate}[label=\textbf{Шаг \arabic*:},leftmargin=*,itemsep=0.35em]
\item Прогнать pm; построить $\Delta|\Phi|$; зафиксировать число пиков и ось.
\item Прогнать pp-контроль на тех же параметрах осреднения.
\item Сопоставить с ожиданием для head-on серии CE-A-02.
\end{enumerate}

\subsubsection*{Анализ, расчёт погрешности и выводы}
pm даёт двухфронтовую прокси, pp --- нет: макро-след различает физику канала знака, а не только масштаб осреднения.
Это необходимый, но не достаточный шаг к ансамблевой статистике CE-A-02 (там --- вращения и $\bar\varepsilon_\pi$).

Тип~A по ориентации при фиксированном seed; пороги пиков --- тип~B.

\textbf{Вывод.}
Lattice-аннигиляция pm оставляет наблюдаемый дипольный след в $|\Phi|$ на~$T$.
\textbf{Статус записи:} опыт закрыт.
